\documentclass[11pt,letterpaper]{article}
\usepackage[top=0.5in, bottom=0.38in, left=0.7in, right=0.7in]{geometry}
\usepackage{lmodern}
\usepackage{hyperref}
\hypersetup{colorlinks=false, urlcolor=black, pdfborder=0 0 0}
\setlength{\parindent}{0pt}
\setlength{\parskip}{0.55em}
\raggedright
\pagestyle{empty}

\begin{document}

\noindent
{\bfseries Ke-Cheng Hsu (Clio)}\\
Ph.D.\ Candidate, Department of Economics, HKUST\\
(852) 6840\,9753 \quad $|$ \quad \href{mailto:khsu@connect.ust.hk}{khsu@connect.ust.hk}\\
\href{https://cliohsu1997.github.io/}{cliohsu1997.github.io} \quad $|$ \quad \href{https://github.com/cliohsu1997}{github.com/cliohsu1997}

\vspace{0.3em}
\noindent
16 November 2026

\vspace{0.2em}
\noindent
Hiring Team\\
QUANTUM BAY INNOVATIONS LIMITED\\
Re: Assistant Quantitative Analyst / 助理量化分析員

\vspace{0.1em}
Dear Hiring Team,

I am writing to apply for the Assistant Quantitative Analyst position. I am a Ph.D.\ candidate in Economics at HKUST (expected June 2027; available as early as February 2027), a Hong Kong permanent resident with the unrestricted right to work in Hong Kong, a native Mandarin speaker, and professionally fluent in English. The posting asks for someone who can frame analysis problems under guidance, set metrics, obtain and analyze data, interpret results, monitor investment risk, prepare regular risk and performance-attribution reports, and strengthen reporting infrastructure---with strong Python/R skills, clear logic, and a willingness to learn. That matches how I already work.

My quantitative foundation matches the degree and tool requirements. I hold a B.Sc.\ in Mathematics and Economics with a minor in Actuarial Mathematics at HKUST (CGA 3.9/4.3; Dean's List; Academic Achievement Medal, top 1\%), an M.Phil.\ in Economics (CGA 4.1/4.3), and I am finishing a Ph.D.\ in empirical industrial organization. I am proficient in Python, R, and C++ (also MATLAB and Stata), and I build reproducible analysis pipelines on AWS and HPC/SLURM. I learn new methods quickly, take ownership of deliverables, and communicate conclusions clearly in Mandarin and English---the learning ability, responsibility, and communication the posting emphasizes.

The Chinese and English job descriptions map directly to my project habits. Under research guidance I define the analysis frame and metrics, pull and clean large panels, run quantitative models, and write decision-ready conclusions. In my sole-authored job market paper on Taobao apparel return policy, I linked orders, prices, insurance, and returns (\(\sim\)62M RMB revenue; six categories; \(\sim\)20--30\% return rates), estimated structural demand and return responses via GMM with bootstrap checks, and produced scenario rankings for profit and welfare (a moderate restocking-fee menu raises social welfare by about 2 million RMB, roughly a 5\% gain, and seller profit by about 1.2 million RMB relative to status-quo policies). That is the same loop as ``understand the problem \(\rightarrow\) choose indicators \(\rightarrow\) analyze \(\rightarrow\) interpret \(\rightarrow\) conclude.''

Daily monitoring and exception reporting are already part of my workflow. With Professors Kohei Kawaguchi and Kei Kawai I built title-by-market-by-month panels from bookstore POS data (deliveries, sales, returns; tens of billions of records across 1{,}000+ stores), estimated pricing and inventory under demand uncertainty, validated estimators with Monte Carlo, and delivered weekly descriptive statistics, regressions, and figures on a fixed deadline---close to aggregating assigned data daily, watching whether results match expectations, and writing a review pack when they do not. The English posting's risk-monitoring and performance-attribution language also matches earlier interbank-crisis modelling in MATLAB, where I traced how liquidity stress and losses propagate across a financial network.

I am ready to support the investment team and senior management with clear risk and performance reports, help improve reporting infrastructure, and grow into business modelling and quantitative analysis under your guidance. Fresh graduates are welcome on this posting; I am an early-career quantitative candidate with deep analytical practice rather than titled fund-management years. My expected monthly basic salary is HKD\,28,000--30,000 (within the advertised HKD\,25,000--35,000 range). My resume is attached. Thank you for your consideration.

\vspace{0.3em}
Sincerely,\\[0.25em]
Ke-Cheng Hsu (Clio)

\end{document}
